% !TEX program = xelatex
\documentclass[10pt,letterpaper,twocolumn]{article}
\usepackage[top=0.7in, bottom=0.7in, left=0.7in, right=0.7in]{geometry}

\usepackage{fontspec}
\usepackage{unicode-math}
\setmainfont{Latin Modern Roman}[Extension=.otf, UprightFont=lmroman10-regular, BoldFont=lmroman10-bold, ItalicFont=lmroman10-italic, BoldItalicFont=lmroman10-bolditalic]
\setsansfont{Latin Modern Sans}[Extension=.otf, UprightFont=lmsans10-regular, BoldFont=lmsans10-bold, ItalicFont=lmsans10-oblique, BoldItalicFont=lmsans10-boldoblique]
\newfontfamily\acbold{ywft-processing-bold}[Path=../../system/public/type/webfonts/, Extension=.ttf]
\newfontfamily\aclight{ywft-processing-light}[Path=../../system/public/type/webfonts/, Extension=.ttf]
\setmonofont{Latin Modern Mono}[Extension=.otf, UprightFont=lmmono10-regular, ItalicFont=lmmono10-italic, Scale=0.85]

\usepackage{xcolor}
\usepackage{titlesec}
\usepackage{enumitem}
\usepackage{booktabs}
\usepackage{tabularx}
\usepackage{fancyhdr}
\usepackage{hyperref}
\usepackage{xurl}
\usepackage{graphicx}
\usepackage{ragged2e}
\usepackage{microtype}

\definecolor{acpink}{RGB}{180,72,135}
\definecolor{acpurple}{RGB}{120,80,180}
\definecolor{acdark}{RGB}{64,56,74}
\definecolor{acgray}{RGB}{119,119,119}
\definecolor{acgreen}{RGB}{40,150,90}
\definecolor{acred}{RGB}{200,50,60}
\definecolor{acamber}{RGB}{200,140,20}

\hypersetup{colorlinks=true, linkcolor=acpurple, urlcolor=acpurple, citecolor=acpurple, pdfauthor={@jeffrey}, pdftitle={No Paint — refresh report, 2026-07-19}}

\titleformat{\section}{\normalfont\bfseries\normalsize\uppercase}{\thesection.}{0.5em}{}
\titlespacing{\section}{0pt}{1.1em}{0.3em}
\titleformat{\subsection}{\normalfont\bfseries\small}{\thesubsection}{0.5em}{}
\titlespacing{\subsection}{0pt}{0.7em}{0.15em}

\pagestyle{fancy}
\fancyhf{}
\renewcommand{\headrulewidth}{0pt}
\fancyhead[C]{\footnotesize\color{acpink}\textit{No Paint — refresh report, for @jeffrey}}
\fancyfoot[C]{\footnotesize\thepage}

\newcommand{\ac}{\textsc{Aesthetic.Computer}}
\newcommand{\ok}{{\color{acgreen}\textbf{✓ done}}}
\newcommand{\open}{{\color{acamber}\textbf{◌ open}}}
\newcommand{\dead}{{\color{acred}\textbf{✗ dead}}}
\newcommand{\pull}[1]{\par\vspace{0.35em}{\aclight\small\color{acpink}\itshape #1}\par\vspace{0.35em}}

\setlist[itemize]{nosep, leftmargin=1.1em, itemsep=0.12em}
\setlist[enumerate]{nosep, leftmargin=1.3em}
\setlength{\columnsep}{1.7em}

\begin{document}

\twocolumn[{%
  \vspace*{0.2em}
  \begin{center}
    {\aclight\large\color{acgray} the no paint refresh}\\[0.5em]
    {\acbold\Huge The Domain Comes Home}\\[0.6em]
    {\large\itshape Hosting swap, a living archive of 32,184 paintings,
      and the road to the .ART Award}\\[0.7em]
    {\small\color{acgray} @jeffrey \;·\; prompt.ac/@jeffrey \;·\; July 19, 2026}\\[0.3em]
    {\small\color{acgray} \url{https://nopaint.art} \;·\; \texttt{grants/dot-art-award-2026/}}
  \end{center}
  \vspace{1.1em}
}]

\section{Summary}

No Paint --- the collaborative painting game that preceded \ac{} ---
moved from Vercel onto AC's own production monolith (lith) today, with
its preserved 2021 build serving at the root and at
\texttt{/classic}, valid TLS, and a new endless gallery over its full
painting archive at \texttt{/gallery}. Along the way the archive turned
out to be far larger and far more alive than expected, the game's save
server turned out to be dead, and lith's TLS layer taught a lesson
worth writing down. The work anchors the entry for the
\textbf{.ART Award 2026} (\$15{,}000 grand prize, deadline
November~1): \emph{No Paint, 2016--present}, an artwork that is
literally a domain.

\section{What shipped}

\begin{itemize}
  \item \ok\ \textbf{Hosting swap.} The complete 2021 Construct export
    (165 files: runtime, 32 spritesheets, all 130 audio files, fonts,
    icons) was mirrored from the live Vercel deployment into
    \texttt{system/public/nopaint.art/} and now serves from lith. DNS
    A records for the apex and \texttt{www} flipped from Vercel
    (76.76.21.21) to lith (209.38.133.33) via \texttt{doctl}; the zone
    was already in DigitalOcean.
  \item \ok\ \textbf{Source recovery.} The editable Construct project
    (\texttt{No Paint.c3p}, 7.4\,MB) was pulled from Dropbox and its
    SHA-256 matches the provenance hash recorded in the migration
    plan. It is extracted locally --- 45 event sheets, 82 layouts,
    205 object types, 1{,}197 images --- so brush ports need neither
    Dropbox nor Construct.
  \item \ok\ \textbf{Gallery.} \url{https://nopaint.art/gallery} ---
    an endless grid over every painting in the archive, with by-year
    filters and a shuffle view, images served straight from the
    existing \texttt{pix.nopaint.art} CDN, styled in the game's own
    Pixelar face and \texttt{\#ebebeb} grey. Driven by a generated
    1.4\,MB manifest (408\,KB gzipped).
  \item \ok\ \textbf{Valid TLS.} Let's Encrypt production certificate,
    \texttt{CN=nopaint.art}, valid through October 17, plus a separate
    certificate for \texttt{www}, which 301s to the apex.
\end{itemize}

\section{The archive is alive}

The project file suggested a finished 2021 artwork. The bucket says
otherwise: \textbf{32{,}184 paintings, 3.9\,GB}, each 256$\times$288
pixels, uploaded continuously for four years after ``Summer 2021.''

\begin{center}
\begin{tabular}{lr}
\toprule
\textbf{Year} & \textbf{Paintings kept} \\
\midrule
2021 & 15{,}407 \\
2022 & 8{,}016 \\
2023 & 5{,}449 \\
2024 & 2{,}112 \\
2025 & 1{,}200 \\
\bottomrule
\end{tabular}
\end{center}

\pull{People kept pressing Paint for four years after the artwork was
``done.''}

The taper to zero is not fading interest --- it is infrastructure
death. The game saves through \texttt{server.nopaint.art}, a CNAME to
\texttt{glitch.edgeapp.net}, and Glitch no longer exists. \dead\
Saving is broken today; restoring it (a small lith endpoint writing
into the same DigitalOcean Space, plus a CNAME repoint) would bring
the archive back to life. One loss to flag: the paintings carry no
painter-handle metadata --- no PNG text chunks, no S3 metadata --- so
attribution, if it was ever recorded, lived in the dead Glitch app.
Worth checking for a Glitch project export before declaring it lost.

\section{The TLS post-mortem}

The swap's only real fight was TLS, and the loud symptom was a decoy.
Three problems stacked:

\begin{enumerate}
  \item \textbf{DNS timing.} Caddy's first ACME attempts fired before
    the flip propagated, so validation reached the old Vercel IP.
  \item \textbf{Staging demotion.} After those failures Caddy demoted
    the host to Let's Encrypt's \emph{staging} endpoint, whose
    resolvers held the stale IP for up to the record's one-hour TTL,
    and its retry backoff stretched to tens of minutes.
  \item \textbf{The real bug: certificates were obtained but never
    served.} lith's Caddyfile is one big \texttt{:443} catch-all that
    pins its certificate selection to the Cloudflare origin
    certificate. The Caddyfile adapter emitted \emph{only} that
    catch-all TLS connection policy --- so every SNI, including
    \texttt{nopaint.art}, was handed the origin certificate even
    after issuance succeeded. \texttt{www} held a valid certificate
    for hours without serving it. The proof and the fix both came
    from \texttt{caddy adapt}: adding a distinct per-site TLS setting
    (\texttt{protocols tls1.2 tls1.3}) forces an SNI-matched policy
    with managed-certificate selection, ordered ahead of the
    catch-all.
\end{enumerate}

The recipe is documented in the Caddyfile itself and in agent memory:
any future direct-DNS (non-Cloudflare) domain on lith needs the same
per-site TLS block, and \texttt{caddy adapt \textbar\ jq} verifies the
policy order before deploying.

\section{The award frame}

The .ART Award 2026 (\url{https://award.art}) is uncannily aligned:
the submission \emph{is} a .art domain URL. Grand prize \$15{,}000
cash, plus month-long residencies in France and Spain, editorial
coverage, and a premium .art domain; deadline November~1, winners
December~3; no fee, no geographic restriction, multiple entries
allowed. The jury includes Jerry Saltz and Regina Harsanyi (Museum of
the Moving Image), whose specialty is variable-media and software
preservation --- exactly what this project practices: a hash-verified
source recovery, a preserved classic build at a stable URL, and a
living native remake at the same address.

\pull{The entry: No Paint, 2016--present --- an artwork that is a
domain, submitted as the domain itself.}

Materials live in \texttt{grants/dot-art-award-2026/} (research +
draft bio, process narrative, documentation plan), with a tracker row
on the papers deadlines page. An optional second entry --- cancelok,
on the free-as-of-July \texttt{nopad.art} --- gets decided by
mid-October only if it matures.

\section{What's next}

In order, toward November~1:

\begin{enumerate}
  \item \open\ \textbf{Restore saving.} Lith endpoint
    $\rightarrow$ \texttt{pix.nopaint.art} Space; repoint the dead
    \texttt{server} CNAME; regenerate the gallery manifest on save.
    The classic game becomes fully alive for the first time since
    2025.
  \item \open\ \textbf{Native remake at the root.} The proposal loop
    already landed in \ac{} (\texttt{disks/nopaint.mjs}, explicit
    choosing/proposing/committing states plus the original weighted
    picker and a test suite). The root page presents it; the classic
    build stays at \texttt{/classic}.
  \item \open\ \textbf{Port the authored brushes} from the extracted
    event sheets --- Softy, Bubbles, Grid Worm, Walker, Banner, Wafer
    first --- treating each sheet as a behavioral score, not code to
    translate.
  \item \open\ \textbf{The story page} (\texttt{/story}): 2016
    $\rightarrow$ Summer 2021 $\rightarrow$ the quiet four years
    $\rightarrow$ \ac{} $\rightarrow$ the remake.
  \item \open\ \textbf{Submit}: trim bio and narrative to the form,
    cut the 30--60s process video, verify collection claims, decide
    the credit line (solo vs.\ the 2021 collaborators), file well
    before the deadline.
\end{enumerate}

\vspace{1em}
\begin{center}
{\aclight\small\color{acgray} pressed together from the July 19
session {\normalfont\small\color{acgray}·} \ac{}
{\normalfont\small\color{acgray}·} the archive does not end here}
\end{center}

\end{document}
